% Time form: Present tense
\newpage
\section{Red Team: Reconnaissance Scenario}\label{sec:reconnaissance-scenario-implementation}

The reconnaissance scenario is an engineered red team scenario that systematically maps an unknown network.
It chains four drills in sequence: host discovery, port scanning, and service enumeration.
If no interface is specified by the operator, the scenario first discovers all active physical interfaces and, if any lacks an IP address, requests one via \gls{DHCP} before scanning.
The scenario results serve as the reproducible baseline for comparison with the \gls{AI}-driven reconnaissance approach.
The full scenario implementation is provided in Listing~\ref{lst:reconnaissance-scenario-py}.

\subsection{Goal}\label{subsec:reconnaissance-scenario-goal}

The goal of the scenario is to identify all reachable hosts, their open ports and running services in a target network without prior knowledge of the topology.
The drill chain is designed to work generically across different network environments: the operator can either provide a specific interface and subnet or let the scenario discover the available interfaces automatically.

\subsection{Discover Hosts Drill}\label{subsec:reconnaissance-drill-discover-hosts}

The \textit{Discover Hosts} drill is the first active step of the scenario.
For local network segments, it runs \texttt{arp-scan --localnet} on the given interface, which sends \gls{ARP} requests and collects the \gls{IP} address, \gls{MAC} address, and vendor of each responding host.
If an explicit subnet is provided, the drill falls back to \texttt{nmap -sn} instead, which uses \gls{ICMP} and \gls{TCP} probes and therefore works across Layer~3 boundaries.
\gls{MAC} addresses are only available for hosts on the same Layer~2 segment; remote hosts are recorded without a \gls{MAC}.
The full implementation is provided in Listing~\ref{lst:discover-hosts-drill-py}.

\subsubsection{Drill Interface}\label{subsubsec:reconnaissance-port-discovery}

The drill exposes two arguments: a required \texttt{interface} and an optional \texttt{subnet}.
If \texttt{subnet} is omitted, the drill defaults to a local ARP scan; if provided, it switches to the nmap-based subnet scan.
The resource configuration is shown in Listing~\ref{lst:discover-hosts-drill-yaml}.

\subsection{Enumeration}\label{subsec:enumeration}

After host discovery, the scenario runs two further drills that progressively deepen the collected information.

\subsubsection{Port Scan}\label{subsubsec:port-scan}

The \textit{Port Scan} drill receives the \texttt{NetworkDiscoveryResultMap} produced by host discovery and runs \texttt{nmap -sV --open} against each discovered \gls{IP}.
For every open port, the drill records the port number, transport protocol (\gls{TCP}/\gls{UDP}), service name, and detected version string and appends a \texttt{NetworkService} entry to the existing result map.
The full implementation is provided in Listing~\ref{lst:port-scan-drill-py}.

\subsubsection{Service Enumeration}\label{subsubsec:service-enumeration}

The \textit{Enumerate Services} drill takes the enriched result map and runs targeted \gls{NSE} scripts for each discovered service.
Well-known services such as HTTP, DNS, FTP, SMTP, SMB, and LDAP are mapped to a specific set of scripts (e.g.\ \texttt{http-title}, \texttt{dns-zone-transfer}, \texttt{smb-security-mode}).
Unknown services fall back to the generic \texttt{banner} script, which captures the initial response of any \gls{TCP} service.
The drill returns a mapping of \texttt{ip:port} keys to lists of finding strings extracted from the \gls{NSE} output.
The full implementation is provided in Listing~\ref{lst:enumerate-services-drill-py}.

\subsection{Output Artifacts}\label{subsec:reconnaissance-scenario-artifacts}

The primary output of the scenario is the \texttt{NetworkDiscoveryResultMap}, which aggregates all results per interface.
Each entry maps an interface name to a \texttt{NetworkDiscoveryResult} containing the list of discovered \texttt{NetworkService} objects with their endpoints (IP, MAC, port, protocol, version).
After port scanning, the scenario prints the result as a human-readable table.
The service enumeration findings are printed separately as a structured list grouped by \texttt{ip:port}.
